\subsubsection{Przegląd ekranów aplikacji}

Interfejs użytkownika znajduje się w folderze \texttt{ui}, a jego struktura przedstawia się
następująco:

\begin{verbatim}
ui/
├── MainScreen.kt
├── components/
│   └── ConnectionStatusBar.kt
├── map/
│   ├── MapScreen.kt
│   └── MapViewModel.kt
├── nodes/
│   ├── NodeDetailScreen.kt
│   ├── NodeFilterState.kt
│   ├── NodeItem.kt
│   └── NodeListScreen.kt
├── settings/
│   ├── SettingsScreen.kt
│   └── SettingsViewModel.kt
├── theme/
│   ├── Color.kt
│   ├── Theme.kt
│   └── Type.kt
└── zones/
    ├── ZoneBottomSheet.kt
    ├── ZoneConfirmDialog.kt
    ├── ZoneDetailScreen.kt
    └── ZoneViewModel.kt
\end{verbatim}

W dolnej części ekranu znajduje się pasek nawigacyjny z trzema zakładkami:
\textbf{Mapy}, \textbf{Węzły} i \textbf{Ustawienia}.

\paragraph{Mapy.}

\texttt{MapScreen} jest głównym ekranem aplikacji. Wyświetla mapę Google Maps,
na której prezentowane są markery węzłów oraz pozycja użytkownika. Każdy marker renderowany
jest jako bitmapa z inicjałami węzła. Kolor markera zależy od roli węzła oraz stanu połączenia.
W prawym górnym rogu markera znajduje się znacznik jakości sygnału. W lewym dolnym rogu ekranu
umieszczono \textit{floating button} otwierający menu geofencingu.

\begin{figure}[!h]
    \caption{Ekran mapy aplikacji MeshTracker}
    \label{fig:map-screen}
    \centering
    \includegraphics[width=0.45\linewidth]{mapaZblizenie.jpg}
\end{figure}

\paragraph{Markery.}

Wizualizacja markerów została zaprojektowana tak, aby w jak najbardziej przejrzysty sposób
prezentować możliwie wiele danych jednocześnie. Kolor markera odpowiada przypisanej mu roli,
a dodatkowa warstwa kolorów reprezentuje stan połączenia węzła z siecią Meshtastic
(rysunek~\ref{fig:markery-kolory}).

\begin{figure}[!h]
    \centering
    \begin{subfigure}{0.48\linewidth}
        \centering
        \includegraphics[width=\linewidth]{markeryRole.png}
        \caption{Mapowanie roli węzła do koloru}
        \label{fig:markery-role}
    \end{subfigure}
    \hfill
    \begin{subfigure}{0.48\linewidth}
        \centering
        \includegraphics[width=\linewidth]{markeryStanPolaczenia.png}
        \caption{Mapowanie stanu połączenia do koloru}
        \label{fig:markery-stan-polaczenia}
    \end{subfigure}
    \caption{Kodowanie kolorystyczne markerów}
    \label{fig:markery-kolory}
\end{figure}

Węzeł lokalny - czyli ten połączony z telefonem użytkownika bezpośrednio przez Bluetooth -
ma małą ikonkę telefonu w prawym górnym rogu. Pozostałe węzły mają w tym miejscu kropkę,
której kolor odpowiada jakości połączenia z siecią (zakres SNR - stosunek mocy użytecznego
sygnału do mocy szumu, wyrażony w~dB). Strzałka wewnątrz markera wskazuje kierunek
poruszania się węzła (rysunek~\ref{fig:markery-szczegoly}).

\begin{figure}[!h]
    \centering
    \begin{subfigure}{0.48\linewidth}
        \centering
        \includegraphics[width=\linewidth]{markeryJakoscPolaczenia.png}
        \caption{Jakość połączenia (SNR) a kolor kropki}
        \label{fig:markery-jakosc-polaczenia}
    \end{subfigure}
    \hfill
    \begin{subfigure}{0.48\linewidth}
        \centering
        
\includegraphics[width=\linewidth]{markeryKierunek.png}
        \caption{Kierunek poruszania się węzła}
        \label{fig:markery-kierunek}
    \end{subfigure}
    \caption{Szczegóły wizualizacji markerów}
    \label{fig:markery-szczegoly}
\end{figure}

\paragraph{Geofencing.}

Menu stref dostępne jest pod \textit{floating button} w lewym dolnym rogu ekranu mapy.
Strefy składają się z wielokątów, których wierzchołki tworzone są poprzez długie przyciśnięcie
(\textit{long press}) wybranego punktu na mapie. Po dodaniu trzeciego wierzchołka pojawia się
opcja \textbf{Zamknij}, która tworzy strefę z wprowadzonych punktów
(rysunek~\ref{fig:geofencing-tworzenie}).

\begin{figure}[!h]
    \centering
    \begin{subfigure}{0.45\linewidth}
        \centering
        \includegraphics[width=\linewidth]{strefyGeofencing.jpg}
        \caption{Menu geofencingu - widok pusty}
        \label{fig:geofencing-puste}
    \end{subfigure}
    \hfill
    \begin{subfigure}{0.45\linewidth}
        \centering
        \includegraphics[width=\linewidth]{strefyGeofencingWierzcholki.jpg}
        \caption{Definiowanie wierzchołków strefy}
        \label{fig:geofencing-wierzcholki}
    \end{subfigure}
    \caption{Tworzenie strefy geofencingu}
    \label{fig:geofencing-tworzenie}
\end{figure}

Następnie użytkownik wprowadza nazwę strefy, nadaje jej kolor oraz przypisuje węzły, które
będą wyzwalać zdarzenia wejścia i wyjścia. Dostępne strefy mogą być aktywowane i usuwane
z poziomu listy (rysunek~\ref{fig:geofencing-zarzadzanie}).

\begin{figure}[!h]
    \centering
    \begin{subfigure}{0.45\linewidth}
        \centering
        \includegraphics[width=\linewidth]{nowaStrefa.jpg}
        \caption{Formularz tworzenia nowej strefy}
        \label{fig:nowa-strefa}
    \end{subfigure}
    \hfill
    \begin{subfigure}{0.45\linewidth}
        \centering
        \includegraphics[width=\linewidth]{strefyLista.jpg}
        \caption{Lista zdefiniowanych stref}
        \label{fig:strefy-lista}
    \end{subfigure}
    \caption{Zarządzanie strefami geofencingu}
    \label{fig:geofencing-zarzadzanie}
\end{figure}
